\section{The critical abelian cost on the original common window}
\label{sec:critical-abelian-actual-window}

\paragraph{Original triples and the retained omitted set.}
Keep the actual host, finite datum, page, prescribed pair and every
original coefficient completion of AC-G190-window.  As there,
$h\geq2$ and $T\geq1$; a larger initial source threshold
ensures this range.  Thus
\[
 \begin{gathered}
 H_u=\mathbb Q(\zeta_{2^{h+2}},\zeta_7,
       \sqrt{-q},\sqrt{-d},\omega,\sqrt[3]{269}),
       \quad K=\mathbb Q(\omega),\\
 u^{2+8\varepsilon}\leq2^h\leq2u^{2+8\varepsilon},
 \quad0<\varepsilon<1/256,\quad T=\lfloor u^{399}\rfloor,
 \quad D_u=144\cdot2^h.
 \end{gathered}\tag{CAW.1}
\]
Let $Z_u$ be exactly $Z^{\rm all7}_2(u)$ of (Q7W.2), and let
$J_u$ be its complement in the original triples
$(\ell,\eta,e)$ of (TC.1), with $1\leq e<T$.
Partition $J_u=J_{u,0}\sqcup J_{u,\rm mix}$ by whether the
Kummer cubic factor of $\eta$ is trivial or nontrivial.
No old omitted triple is restored.

For clarity distinguish the degree of the abelian rational field
from $D_u$:
\[
 \begin{gathered}
 F_0=\mathbb Q(\zeta_{2^{h+2}},\zeta_7,\sqrt{-q},\sqrt{-d}),
       \qquad F=F_0K,\qquad d_u=[F:\mathbb Q]=96\cdot2^h,\\
 \Delta_u=|\operatorname{disc}F|,
 \qquad f_u=84\cdot2^hqd,
 \qquad R_u^{\rm rad}=2\cdot3\cdot7\cdot269\cdot qd.
 \end{gathered}\tag{CAW.2}
\]
For $(\ell,\eta,e)\in J_{u,0}$, let
$\psi_{\eta,e}^{\rm cr}$ and $\psi_{\eta,e}^{\rm nc}$ be its
two rational finite-character extensions in the same original
$O_\ell$, with complex-conjugation values
$(-1)^{e+1}$ and $(-1)^e$, respectively.
Use exactly the rational strict groups of AC-G191: full local
zero at $\ell$, inertia zero at other finite localization primes,
and the real-place condition, including at coefficient two.
Their Tate twist is still $-e$.

\paragraph{Two zero noncritical rows at the original coefficient three.}
For $h\geq3$, let $V_u\subset J_{u,0}$ consist of
\[
 (3,\delta|_{G_K},e),\qquad 1\leq e<T,
 \quad \delta\text{ even primitive of conductor }32,
 \quad\overline{\delta(5)}^{\,2}
             -\overline{\delta(5)}-1=0.
 \tag{CAW.3}
\]
The reduction uses the same prescribed original $O_3$.
For $h=2$, put $V_u=\varnothing$.
By AC-G193 their noncritical rational strict groups are zero.
Define the explicit nonnegative arithmetic remainder
\[
 \begin{split}
 \mathcal R_{399}^{\rm nc}(u)={}&
 \sum_{(\ell,\eta,e)\in J_{u,\rm mix}}
                \frac{\log\#\mathcal S_{\ell,\eta}(e)}{b_\ell}\\
 &+\sum_{(\ell,\eta,e)\in J_{u,0}\setminus V_u}
       \frac{\log\#\mathcal S^{\mathbb Q}_{\ell,
                       \psi_{\eta,e}^{\rm nc}}(e)}{b_\ell}.
 \end{split}\tag{CAW.4}
\]
It keeps every retained nontrivial cubic strict group and every
other retained noncritical rational component.  This is not a
decomposition of positive derived indices.

\begin{theorem}[AC-G192-window: a critical-component remainder on the same window]
\label{AC-G192-window}
The quantities in (CAW.4) are finite and use the original
coefficient primes, orders, embeddings and full local conditions.
Put
\[
 \begin{split}
 \mathcal P_u(T)={}&d_uT^2\log(2d_uTf_uR_u^{\rm rad})\\
 &+Td_u^4\log(2d_u\Delta_uR_u^{\rm rad}),\\
 \mathcal E_u(T)={}&8D_u^2T^2\log(8TD_u)
       +\tfrac12D_u\mathcal R_uT^2+2D_ur_uc_0T,\\
 \mathcal L_u(T)={}&8D_u^2T^2\log(8TD_u),\\
 \mathcal A_u(T)={}&C_1\,2^hT^2\bigl(h+\log(T+2)+1\bigr),
 \end{split}\tag{CAW.5}
\]
where $\mathcal R_u,r_u,c_0$ are the original (FL.1) data,
and $C_1$ is an absolute constant from (G189.4).
There is an absolute $C_2$ such that
\[
 \begin{aligned}
 \bigl|I_{399}-\mathcal R_{399}^{\rm nc}(u)\bigr|
   &\leq \mathcal E_u(T)+\mathcal L_u(T)
                   +\mathcal A_u(T)+C_2\mathcal P_u(T),\\
 \bigl|I_{399}^{\rm all7,2}-\mathcal R_{399}^{\rm nc}(u)\bigr|
   &\leq \mathcal E_u(T)+\mathcal L_u(T)+C_2\mathcal P_u(T).
 \end{aligned}\tag{CAW.6}
\]
The first is a single full-union local error, one degree-one
torsion bound, the already proved quadratic-seven strict cost,
and one aggregate critical-component cost.
Their combined size is
\[
 O_\varepsilon\bigl(u^{806+16\varepsilon}(\log u)^2\bigr)
                 =o\bigl(\exp(u^\varepsilon/2)\bigr),
 \tag{CAW.7}
\]
uniformly in the original prescribed pair and page.
No character, coefficient-prime or family-count multiplier is added.

Consequently a bound for (CAW.4) with the exact (TC.3) quantifiers
and positive relative slack supplies the full original numerical
premise on the identical $W$.  Conversely the preceding
$I_{399}^{\rm all7,2}$ premise supplies this new remainder premise
on its same $W$, after a threshold increase and slack reduction.
This converse uses the comparison (CAW.6), not a claimed pointwise
ordering between the two different sums.

The set $V_u$ has cardinality
$2\mathbf1_{h\geq3}(T-1)$.  Its removal in (CAW.4) is an exact
zero in the noncritical rational remainder.  It does not assert
zero for the original $K$-strict groups or their positive indices.
Every pure two-power coefficient-three label is now accounted for
either by the old strict-zero union or by this critical-cost
comparison and AC-G193.  No new conclusion about the mixed or
auxiliary remainder is asserted.
\end{theorem}

\begin{proof}
\proofstep{1. The original rational dictionary and its parity.}
For a trivial-cubic character the finite extension it defines
lies in $F/K$.  Conductor disjointness gives $F_0\cap K=\mathbb Q$.
Restriction of characters from $\operatorname{Gal}(F/\mathbb Q)$
to $\operatorname{Gal}(F/K)$ has kernel $\{1,\eta_K\}$,
with $\eta_K(c)=-1$.
There are therefore exactly two extensions and exactly one
of each specified parity.  The rational group has exponent
$3\cdot2^h$, already used in (FC.1), so their values belong
to the very same original $O_\ell$.
No coefficient field, lattice or residue embedding is replaced.

AC-G191 compares these extensions integrally.  At odd primes
the index-two induced lattice splits on the original ring;
at two its determinant-two cokernel and the real-place term
are retained in the proved finite error.
In particular no integral division by two is made at coefficient
two, and the real-place condition in (CAW.4) is not dropped.
The argument applies to every retained subset, without assuming
that $J_u$ or its original-residue exceptions are Galois-stable.
It includes the characters with the growing $q,d$ factors.

\proofstep{2. The exact two coefficient-three rows.}
At $h\geq3$ the original host contains $\mathbb Q(\zeta_{32})$.
There are four even primitive conductor-$32$ characters:
their values on $5$ are the primitive eighth roots of unity.
In characteristic three,
\[
 X^4+1=(X^2+X-1)(X^2-X-1),
 \tag{CAW.8}
\]
with distinct factors and two roots each in the original residue
field, which contains the eighth roots when $h\geq3$.
The original embedding fixes which two labels satisfy each factor.
The two with the plus-linear factor were in G180 and hence $Z_u$;
the two with the minus-linear factor are exactly (CAW.3).
The complete coefficient-three accounting in G188 shows that
these and only these pure two-power labels remained in $J_u$.
Distinct two-power characters remain distinct on $G_K$, because
their ratio cannot be the conductor-three character $\eta_K$.
This proves $|V_u|=2(T-1)$ when $h\geq3$ and the stated zero
count when $h=2$.

For an even $\delta$ in (CAW.3), the noncritical rational finite
extension is exactly $\delta\eta_K^e$.
AC-G193 applies on its unchanged $O_3$ and gives
\[
       \mathcal S^{\mathbb Q}_{3,\psi_{\delta,e}^{\rm nc}}(e)=0.
 \tag{CAW.9}
\]
Its full-local definition agrees with AC-G191: outside the finite
localization set, inflation already kills inertia, and the
real-place condition is automatic at coefficient three.
The exact index-two identity (N32.4) identifies the original
$K$-strict group with its critical rational component; it does
not make that component zero.
Thus the sum in (CAW.4) equals the unreduced $\mathcal R_{J_u}$
of (CP.3) exactly, with no error or new coefficient normalization.

\proofstep{3. One aggregate bound for the retained strict sum.}
Finiteness of these strict sums and their rational terms is part
of AC-G191; the original index and torsion sums are finite by
AC-K27 and (LT.2).  Thus no difference of infinite sums is used.
Write $\mathcal B_{J_u}$ for the normalized original strict sum
over $J_u$, and $\mathcal B_{Z_u}$ for the sum over the old
omitted set $Z_u$.  AC-G191, applied literally with this $J_u$,
gives
\[
       |\mathcal B_{J_u}-\mathcal R_{399}^{\rm nc}|
                            \leq C_2\mathcal P_u(T).
 \tag{CAW.10}
\]
This is the strict-sum comparison (CP.4); its separate
positive-index comparison (CP.24) is not added again below.
All retained critical terms form one nonnegative subsum of the
aggregate all-prime critical sum (CP.2).
That bound covers the growing abelian labels, all positive
weights below $T$, and both coefficient primes dividing the
character orders with their integral errors.

Every strict group in $Z_u$ is zero except the bounded
quadratic-seven coefficient-two family added by G190.
The literal estimate (G189.4) therefore gives
\[
                 0\leq\mathcal B_{Z_u}\leq\mathcal A_u(T).
 \tag{CAW.11}
\]
The strict-group sum over all original triples is
$\mathcal B=\mathcal B_{J_u}+\mathcal B_{Z_u}$.
Consequently (CAW.10)--(CAW.11) compare that whole sum with
(CAW.4) at cost at most $C_2\mathcal P_u(T)+\mathcal A_u(T)$.
No second critical bound is added for the two rows in $V_u$:
their critical components were already in the same aggregate
bound over $J_{u,0}$.

\proofstep{4. One full local error and the original torsion term.}
Let $\mathcal H$ be the normalized sum of the original
$\log\#H^2(C_{\ell,\eta,e})$ over all triples.
The absolute error (FL.4), summed once over $1\leq e<T$,
gives $|\mathcal H-\mathcal B|\leq\mathcal E_u(T)$.
This one estimate covers $Z_u$ and $J_u$ together.

The determinant inequality (EU.2) gives
\[
 0\leq\mathcal H-I_{399}
 \leq\sum_{1\leq e<T}\sum_\ell\frac1{b_\ell}
      \sum_\eta\log\#H^1(C_{\ell,\eta,e})_{\rm tors}.
 \tag{CAW.12}
\]
By (LT.2), the summand at weight $e$ is at most
$8(e+1)D_u^2\log(8(e+1)D_u)$.
For $T\geq1$, $\sum_{1\leq e<T}(e+1)\leq T^2$,
so (CAW.12) is at most $\mathcal L_u(T)$.
This retains the original degree-one torsion and the full
normalization $b_\ell$, including coefficient ramification.
Combining these two errors with Step 3 proves the first line
of (CAW.6).

For its second line apply the same termwise local and torsion
inequalities just to $J_u$; their nonnegative sums are bounded
by the same complete $\mathcal E_u(T)$ and $\mathcal L_u(T)$.
Then use (CAW.10), with no $\mathcal B_{Z_u}$ term.
The two displayed comparisons are separate conclusions; their
full local errors are not summed with one another.
At no point is a positive derived index split into two rational
indices.  The comparison passes through the actual finite
strict groups and the actual integral determinant.

\proofstep{5. Uniform source growth and the identical window.}
AC-G191's (CP.5) and (CP.23) give
$\mathcal P_u(T)=O_\varepsilon(u^{806+16\varepsilon}(\log u)^2)$
for the actual growing $F,q,d,h$.
The distinction $d_u=96\cdot2^h$ and $D_u=144\cdot2^h$
is retained in (CAW.5).
The old local and torsion terms are
$O_\varepsilon(u^{802+16\varepsilon}(\log u)^2)$,
and the quadratic-seven cost is
$O_\varepsilon(u^{800+8\varepsilon}\log u)$.
These prove (CAW.7), uniformly over the prescribed source pairs.

Fix a retained finite datum, page and $0<\varepsilon<1/256$.
Assume the exact (TC.3) premise with (CAW.4) in place of the
full index sum.  Thus $J\geq1$, $X_0$ and $s\in(0,1)$ are
independent of $X,q,d,u,\ell,\eta,e$, and each prescribed pair
at $X\geq X_0$ supplies one $W$ with $X^{1/J}\leq W$,
$W^2\leq X$ and
$\mathcal R_{399}^{\rm nc}(u)\leq(1-s)\exp(u^\varepsilon/2)$
for every $u\in[W,W^2]$.
Choose $U_0$ so that the error in the first line of (CAW.6)
is at most $(s/2)\exp(u^\varepsilon/2)$ for every permitted
$u\geq U_0$.
Raise the threshold to $\max(X_0,U_0^J)$ and keep that exact $W$.
It now lies above $U_0$, and the full numerical premise holds
on the same interval with slack $s/2$.

Starting instead with the previous $I_{399}^{\rm all7,2}$
premise, the second line of (CAW.6) gives the new premise by
the same threshold and half-slack argument.
Uniform little-$o$ formulations and fixed finite simultaneous
families use one maximum of their uniform error constants.
Separate existential windows are never intersected.
Every host, page, prescribed pair, original coefficient and actual
source rank $R.t$ is unchanged through the retained K35 interface.
The finite Fourier comparison costs from (FC.4)--(FC.5) remain
part of that original interface and are not replaced by an
integral character projector.
\end{proof}

\paragraph{What remains.}
The explicit sum (CAW.4) still contains all retained nontrivial
cubic sectors and all unsupported noncritical abelian sectors,
including growing auxiliary characters.
Its numerical upper bound remains unproved.
The mixed theta/Selmer realization $\mathcal R_{\rm cub}$ and
H04.4 remain separate open inputs with their original normalizations.
No stronger AW.1 source package, unproved window compatibility,
coefficient replacement or rational regulator-to-lattice inference
has been introduced.  The retained finite-cover/Artin--Tate
dictionary is unchanged.  This section composes the fully checked
critical bound and noncritical zero with the actual arithmetic
gate; it does not reprove their deep special-value inputs or
claim Lean execution.
